\begin{longtable}{>{\hangindent.5cm\hangafter=1\raggedright}p{3cm}
                    >{\hangindent.5cm\hangafter=1\raggedright}p{3cm}
                    >{\hangindent.5cm\hangafter=1\raggedright}p{3cm}
                    >{\hangindent.5cm\hangafter=1\raggedright}p{3cm}}
\caption{Global Attributes used in ECCO V4r4 data netCDF files}
\label{tab:variable-attributes} \\ 
\hline \endfirsthead
\endhead
\endfoot
\hline \endlastfoot
\rowcolor{blue!25} \textbf{Attribute Name} & \textbf{Format} & \textbf{Description} & \textbf{Source} \\ \hline
\rowcolor{cyan!25}
conventions & string & A comma-separated list of the conventions that are followed by the dataset. For files that follow this version of ACDD, include the string 'ACDD-1.3'. (This attribute is described in the NetCDF Users Guide.) & ACDD \\ \hline
\rowcolor{cyan!25}
acknowledgement & string & A place to acknowledge various types of support for the project that produced this data. & ACDD \\ \hline
\rowcolor{cyan!25}
naming\_authority & string & 	The organization that provides the initial id (see above) for the dataset. The naming authority should be uniquely specified by this attribute. We recommend using reverse-DNS naming for the naming authority; URIs are also acceptable. Example: 'edu.ucar.unidata'. & ACDD \\ \hline
\rowcolor{cyan!25}
cdm\_data\_type & string & The data type, as derived from Unidata's Common Data Model Scientific Data types and understood by THREDDS. This is a THREDDS 'dataType', and is different from the CF NetCDF attribute 'featureType', which indicates a Discrete Sampling Geometry file in CF. For ECCO, it is 'grid'. & ACDD \\ \hline
\rowcolor{cyan!25}
comment & string & Miscellaneous information about the data, not captured elsewhere. This attribute is defined in the CF Conventions. & CF, ACDD \\ \hline
\rowcolor{cyan!25}
coordinates & string & This is necessary if the variable has coordinates that are not the same as the dimension names of the data. In this case, the coordinates attribute identifies the auxiliary variables that contain geospatial or temporal coordinates. For example,  variables on a trajectory often use the data acquisition time as a single dimension but have auxiliary coordinates that record the latitude and longitude for each datum in the main variable. If a variable “foo” has a coordinate of time but co- exists with variables named “lat” and “lon” with the corresponding geolocation information, then the coordinates variable would read “lat lon” to indicate that the variables hold the appropriate geospatial location information. Without this attribute, tools will not be able to locate the corresponding latitude and longitude for data in a swath, trajectory, or non-rectilinear grid. Note that the auxiliary coordinate variables must still follow the conventions for units that are noted in Section 4.4 and Section 4.5 above. Also, note that the only delimiter between words inside the coordinates attribute should be a space. & CF \\ \hline
\rowcolor{cyan!25}
creator\_email & string & The email address of the person (or other creator type specified by the creator\_type attribute) principally responsible for creating this data. & ACDD \\ \hline
\rowcolor{cyan!25}
institution & string & The name of the institution principally responsible for originating this data. This attribute is recommended by the CF convention. & ACDD \\ \hline
\rowcolor{cyan!25}
creator\_name & string & The name of the person (or other creator type specified by the creator\_type attribute) principally responsible for creating this data. & ACDD \\ \hline
\rowcolor{cyan!25}
creator\_type & string & Specifies type of creator with one of the following: 'person', 'group', 'institution', or 'position'. If this attribute is not specified, the creator is assumed to be a person. & ACDD \\ \hline
\rowcolor{cyan!25}
creator\_url & string & The URL of the person (or other creator type specified by the creator\_type attribute) principally responsible for creating this data. & ACDD \\ \hline
\rowcolor{cyan!25}
date\_created & string & The date on which this version of the data was created. (Modification of values implies a new version, hence this would be assigned the date of the most recent values modification.) Metadata changes are not considered when assigning the date\_created. The ISO 8601:2004 extended date format is recommended, as described in the Attribute Content Guidance section. & ACDD \\ \hline
\rowcolor{cyan!25}
date\_issued & string & The date on which this data (including all modifications) was formally issued (i.e., made available to a wider audience). Note that these apply just to the data, not the metadata. The ISO 8601:2004 extended date format is recommended, as described in the Attributes Content Guidance section. & ACDD \\ \hline
\rowcolor{cyan!25}
date\_metadata\_modified & string & The date on which the metadata was last modified. The ISO 8601:2004 extended date format is recommended, as described in the Attributes Content Guidance section. & ACDD \\ \hline
\rowcolor{cyan!25}
date\_modified & string & The date on which the data was last modified. Note that this applies just to the data, not the metadata. The ISO 8601:2004 extended date format is recommended, as described in the Attributes Content Guidance section. & ACDD \\ \hline
\rowcolor{cyan!25}
geospatial\_bounds\_crs & string & Describes the data's 2D or 3D geospatial extent in OGC's Well-Known Text (WKT) Geometry format (reference the OGC Simple Feature Access (SFA) specification). The meaning and order of values for each point's coordinates depends on the coordinate reference system (CRS). The ACDD default is 2D geometry in the EPSG:4326 coordinate reference system. The default may be overridden with geospatial\_bounds\_crs and geospatial\_bounds\_vertical\_crs (see those attributes). EPSG:4326 coordinate values are latitude (decimal degrees\_north) and longitude (decimal degrees\_east), in that order. Longitude values in the default case are limited to the [-180, 180) range. Example: 'POLYGON ((40.26 -111.29, 41.26 -111.29, 41.26 -110.29, 40.26 -110.29, 40.26 -111.29))'. & ACDD \\ \hline
\rowcolor{cyan!25}
geospatial\_lat\_max & float & Describes a simple upper latitude limit; may be part of a 2- or 3-dimensional bounding region. Geospatial\_lat\_max specifies the northernmost latitude covered by the dataset.& ACDD \\ \hline
\rowcolor{cyan!25}
geospatial\_lat\_min & float & Describes a simple lower latitude limit; may be part of a 2- or 3-dimensional bounding region. Geospatial\_lat\_min specifies the southernmost latitude covered by the dataset. & ACDD \\ \hline
\rowcolor{cyan!25}
geospatial\_lat\_resolution & string & Information about the targeted spacing of points in latitude. Recommend describing resolution as a number value followed by the units. Examples: '100 meters', '0.1 degree' & ACDD \\ \hline
\rowcolor{cyan!25}
geospatial\_lat\_units & string & Units for the latitude axis described in "geospatial\_lat\_min" and "geospatial\_lat\_max" attributes. These are presumed to be "degree\_north"; other options from udunits may be specified instead.& ACDD \\ \hline
\rowcolor{cyan!25}
geospatial\_lon\_max & float & Describes a simple longitude limit; may be part of a 2- or 3-dimensional bounding region. geospatial\_lon\_max specifies the easternmost longitude covered by the dataset. Cases where geospatial\_lon\_min is greater than geospatial\_lon\_max indicate the bounding box extends from geospatial\_lon\_max, through the longitude range discontinuity meridian (either the antimeridian for -180:180 values, or Prime Meridian for 0:360 values), to geospatial\_lon\_min; for example, geospatial\_lon\_min=170 and geospatial\_lon\_max=-175 incorporates 15 degrees of longitude (ranges 170 to 180 and -180 to -175). & ACDD \\ \hline
\rowcolor{cyan!25}
geospatial\_lon\_min & float & Describes a simple longitude limit; may be part of a 2- or 3-dimensional bounding region. geospatial\_lon\_min specifies the westernmost longitude covered by the dataset. See also geospatial\_lon\_max. & ACDD \\ \hline
\rowcolor{cyan!25}
geospatial\_lon\_resolution & float & Information about the targeted spacing of points in longitude. Recommend describing resolution as a number value followed by units. Examples: '100 meters', '0.1 degree' & ACDD \\ \hline
\rowcolor{cyan!25}
geospatial\_lon\_units & string & Units for the longitude axis described in "geospatial\_lon\_min" and "geospatial\_lon\_max" attributes. These are presumed to be "degree\_east"; other options from udunits may be specified instead. & ACDD \\ \hline
\rowcolor{cyan!25}
geospatial\_vertical\_max & float & Describes the numerically larger vertical limit; may be part of a 2- or 3-dimensional bounding region. See geospatial\_vertical\_positive and geospatial\_vertical\_units. & ACDD \\ \hline
\rowcolor{cyan!25}
geospatial\_vertical\_min & float & Describes the numerically smaller vertical limit; may be part of a 2- or 3-dimensional bounding region. See geospatial\_vertical\_positive and geospatial\_vertical\_units. & ACDD \\ \hline
\rowcolor{cyan!25}
geospatial\_vertical\_positive & string & One of 'up' or 'down'. If up, vertical values are interpreted as 'altitude', with negative values corresponding to below the reference datum (e.g., under water). If down, vertical values are interpreted as 'depth', positive values correspond to below the reference datum. Note that if geospatial\_vertical\_positive is down ('depth' orientation), the geospatial\_vertical\_min attribute specifies the data's vertical location furthest from the earth's center, and the geospatial\_vertical\_max attribute specifies the location closest to the earth's center. & ACDD \\ \hline
\rowcolor{cyan!25}
geospatial\_vertical\_resolution & string & Information about the targeted vertical spacing of points. Example: '25 meters' & ACDD \\ \hline
\rowcolor{cyan!25}
geospatial\_vertical\_units & string & Units for the vertical axis described in "geospatial\_vertical\_min" and "geospatial\_vertical\_max" attributes. The default is EPSG:4979 (height above the ellipsoid, in meters); other vertical coordinate reference systems may be specified. Note that the common oceanographic practice of using pressure for a vertical coordinate, while not strictly a depth, can be specified using the unit bar. Examples: 'EPSG:5829' (instantaneous height above sea level), 'EPSG:5831' (instantaneous depth below sea level). & ACDD \\ \hline
\rowcolor{cyan!25}
history & string & Provides an audit trail for modifications to the original data. This attribute is also in the NetCDF Users Guide: 'This is a character array with a line for each invocation of a program that has modified the dataset. Well-behaved generic netCDF applications should append a line containing: date, time of day, user name, program name and command arguments.' To include a more complete description you can append a reference to an ISO Lineage entity; see NOAA EDM ISO Lineage guidance. & CF, ACDD \\ \hline
\rowcolor{cyan!25}
id & string & An identifier for the data set, provided by and unique within its naming authority. The combination of the "naming authority" and the "id" should be globally unique, but the id can be globally unique by itself also. IDs can be URLs, URNs, DOIs, meaningful text strings, a local key, or any other unique string of characters. The id should not include white space characters. & ACDD \\ \hline
\rowcolor{cyan!25}
institution & string & The name of the institution principally responsible for originating this data. This attribute is recommended by the CF convention. & CF, ACDD \\ \hline
\rowcolor{cyan!25}
instrument\_vocabulary & string & Controlled vocabulary for the names used in the "instrument" attribute. & ACDD \\ \hline
\rowcolor{cyan!25}
keywords & string & A comma-separated list of key words and/or phrases. Keywords may be common words or phrases, terms from a controlled vocabulary (GCMD is often used), or URIs for terms from a controlled vocabulary (see also "keywords\_vocabulary" attribute). & ACDD \\ \hline
\rowcolor{cyan!25}
keywords\_vocabulary & string & If you are using a controlled vocabulary for the words/phrases in your "keywords" attribute, this is the unique name or identifier of the vocabulary from which keywords are taken. If more than one keyword vocabulary is used, each may be presented with a prefix and a following comma, so that keywords may optionally be prefixed with the controlled vocabulary key. Example: 'GCMD:GCMD Keywords, CF:NetCDF COARDS Climate and Forecast Standard Names'. & ACDD \\ \hline
\rowcolor{cyan!25}
license & string & Provide the URL to a standard or specific license, enter 'Public Domain', 'Freely Distributed' or 'None', or describe any restrictions to data access and distribution in free text & ACDD \\ \hline
\rowcolor{cyan!25}
metadata\_link & string & A URL that gives the location of more complete metadata. A persistent URL is recommended for this attribute. & ACDD \\ \hline
\rowcolor{cyan!25}
naming\_authority & string & The organization that provides the initial id (see above) for the dataset. The naming authority should be uniquely specified by this attribute. We recommend using reverse-DNS naming for the naming authority; URIs are also acceptable. Example: 'edu.ucar.unidata'. & ACDD \\ \hline
\rowcolor{cyan!25}
platform & string & Name of the platform(s) that supported the sensor data used to create this data set or product. Platforms can be of any type, including satellite, ship, station, aircraft or other. Indicate controlled vocabulary used in platform\_vocabulary. & ACDD \\ \hline
\rowcolor{cyan!25}
platform\_vocabulary & string & Controlled vocabulary for the names used in the "platform" attribute. & ACDD \\ \hline
\rowcolor{cyan!25}
processing\_level & string & A textual description of the processing (or quality control) level of the data. For ECCO, it is level 4 (L4) & ACDD, GDS \\ \hline
\rowcolor{cyan!25}
product\_version & string & Version identifier of the data file or product as assigned by the data creator. For example, a new algorithm or methodology could result in a new product\_version. & GDS \\ \hline
\rowcolor{cyan!25}
program & string & The overarching program(s) of which the dataset is a part. A program consists of a set (or portfolio) of related and possibly interdependent projects that meet an overarching objective. Examples: 'GHRSST', 'NOAA CDR', 'NASA EOS', 'JPSS', 'GOES-R'. & ACDD \\ \hline
\rowcolor{cyan!25}
project & string & The name of the project(s) principally responsible for originating this data. Multiple projects can be separated by commas, as described under Attribute Content Guidelines. Examples: 'PATMOS-X', 'Extended Continental Shelf Project'. & ACDD \\ \hline
\rowcolor{cyan!25}
publisher\_email & string & The email address of the person (or other entity specified by the publisher\_type attribute) responsible for publishing the data file or product to users, with its current metadata and format. For ECCO, it is: podaac@podaac.jpl.nasa.gov & ACDD \\ \hline
\rowcolor{cyan!25}
publisher\_institution & string & The institution that presented the data file or equivalent product to users; should uniquely identify the institution. If publisher\_type is institution, this should have the same value as publisher\_name. & ACDD \\ \hline
\rowcolor{cyan!25}
publisher\_name & string & The name of the person (or other entity specified by the publisher\_type attribute) responsible for publishing the data file or product to users, with its current metadata and format. & ACDD \\ \hline
\rowcolor{cyan!25}
publisher\_type & string & Specifies type of publisher with one of the following: 'person', 'group', 'institution', or 'position'. If this attribute is not specified, the publisher is assumed to be a person. & ACDD \\ \hline
\rowcolor{cyan!25}
publisher\_url & string & The URL of the person (or other entity specified by the publisher\_type attribute) responsible for publishing the data file or product to users, with its current metadata and format. For ECCO, it is : https://podaac.jpl.nasa.gov & ACDD \\ \hline
\rowcolor{cyan!25}
references & string & Published or web-based references that describe the data or methods used to produce it. Recommend URIs (such as a URL or DOI) for papers or other references. This attribute is defined in the CF conventions. & CF \\ \hline
\rowcolor{cyan!25}
source & string & The method of production of the original data. If it was model-generated, source should name the model and its version. If it is observational, source should characterize it. This attribute is defined in the CF Conventions. Examples: 'temperature from CTD \#1234'; 'world model v.0.1'. & CF \\ \hline
\rowcolor{cyan!25}
standard\_name\_vocabulary & string & The name and version of the controlled vocabulary from which variable standard names are taken. (Values for any standard\_name attribute must come from the CF Standard Names vocabulary for the data file or product to comply with CF.) Example: 'CF Standard Name Table v27'. & ACDD \\ \hline
\rowcolor{cyan!25}
summary & string & A paragraph describing the dataset, analogous to an abstract for a paper. & ACDD \\ \hline
\rowcolor{cyan!25}
time\_coverage\_end & string & Describes the time of the last data point in the data set. Use ISO 8601:2004 date format, preferably the extended format as recommended in the Attribute Content Guidance section. & ACDD \\ \hline
\rowcolor{cyan!25}
time\_coverage\_start & string & Describes the time of the first data point in the data set. Use the ISO 8601:2004 date format, preferably the extended format as recommended in the Attribute Content Guidance section. & ACDD \\ \hline
\rowcolor{cyan!25}
time\_coverage\_duration & string & Describes the duration of the data set. Use ISO 8601:2004 duration format, preferably the extended format as recommended in the Attribute Content Guidance section. & ACDD \\ \hline
\rowcolor{cyan!25}
time\_coverage\_resolution & string & Describes the targeted time period between each value in the data set. Use ISO 8601:2004 duration format, preferably the extended format as recommended in the Attribute Content Guidance section. & ACDD \\ \hline
\rowcolor{cyan!25}
title & string & A short phrase or sentence describing the dataset. In many discovery systems, the title will be displayed in the results list from a search, and therefore should be human readable and reasonable to display in a list of such names. This attribute is also recommended by the NetCDF Users Guide and the CF conventions. & CF, ACDD \\ \hline
\end{longtable}
